\begin{theorem}[Source result P3K.113 under continuation review]\label{night:P3K.113}
\StatementRecord{night:P3K.113}{open}{}{\cref{night:P3K.96}, \cref{night:P3K.109}, \cref{night:P3K.110}; the original source, point line, roots, characters, reference, selected fan and scalar range stated here}
Retain the original marked source
\[
 E_r:\quad y^2=x^3+rx^2+tx+t^2,\qquad G=(0,t),\qquad
 r\in\mathcal V=\mathbb C\setminus\{0,\tfrac14,\tfrac13\}.
\]
Choose either simple root \(t_c\) of
\(27t^2+(4-18r)t+r^2(4r-1)=0\), with its original selected
height \(q=t-t_c\), branch marking and prescribed lattice-basis
fan. Use the genuine point line \(M_a=\mathcal O([aG]-[O])\),
translation by \(-bG\), and either full cap progression
\[
 (a,b;m)=(2,3;3l),\quad l\ge1,
 \qquad (a,b;m)=(6,1;3l+2),\quad l\ge0.
\]
Retain either primitive norm \(n=\pm1\), every prescribed finite
character \(\zeta_m^d\), the full annular reference of P3K.96,
the ordinary attachments at the other finite fibres, and a fixed
nonzero scalar in the strict contraction range. Let \(X\) be the
resulting completed complex threefold.

Every holomorphic isomorphism between two completed scalar endpoints
of this fixed marked source preserves the original base pointwise.
On each regular fibre its lift carries the actual cyclic cover
\(M_a^\times\) and its physical fibres to the corresponding
target cover and physical fibres. The induced source map is a
degree-one affine holomorphic isomorphism of the original elliptic
fibres. The complete proof includes its affine inverse and the full
integral kernel lattice.

Multiplication on the genuine point line, carried through all root,
reference and fan charts, gives the exact group isomorphism
\[
 \mathbb C^*_{\rm physical}
       \xrightarrow{\ \cong\ }\operatorname{Aut}_{\rm hol}(X).
\]
Every such automorphism is isotopic to the identity through this
physical action. The ratio of two endpoint identifications at the
same scalar winding is exactly a physical scalar, with its completed
isotopy and inverse retained.

On the scalar cover keep
\[
 \lambda_\theta=e^{2\pi i\theta}\lambda_*,\qquad
 R_\theta=e^{2\pi i\theta/m}R_*,\qquad R_*^m=\lambda_*.
\]
With the full SEF reference restoration, every holomorphic endpoint
has equal original scalars at its two ends and integral winding
\(J\). All endpoint linear parts on the ordinary torus are \(+I\).
The complete groups of possible windings are
\[
 \begin{array}{c|c}
 (a,b;m)&\text{scalar winding group}\\\hline
 (2,3;3l)&l\mathbb Z\\
 (6,1;3l+2)&m\mathbb Z.
 \end{array}
\]
For M2 put \(s=1/4-r\) and retain the full source lift
\[
 f_2(p)=\frac{y-t-x/2}{x-s},\qquad
 B_{2,1}(p)=f_2(p)f_2(p-G),\qquad U_G=B_{2,1}^{-1},\qquad S=tU_G.
\]
The inverse map evaluates its coefficient at the actual translated
source point. On the good root retain
\(U_{G,\rm old}=z^{2l}U'_G\). The primitive forward endpoint
at \(J=l\) is
\[
 E_R^{\rm ord}=(R^{-l}S)^n,\qquad
 E_R^{\rm cap}=(R^{-l}U'_G)^n.
\]
Their three chronological composites, at
\(R,\zeta_3R,\zeta_3^2R\), are exactly
\[
 E_2^{\rm ord}E_1^{\rm ord}E_0^{\rm ord}=q^nB_\lambda^{-n},
 \qquad
 E_2^{\rm cap}E_1^{\rm cap}E_0^{\rm cap}=h^nB_\lambda^{-n},
 \quad h=q_\infty q=1-t_cz^m.
\]
For M6 the primitive winding \(J=m\) is the canonical forward
\(q^n/h^n\) quotient closing. Negative windings use the exact
chronological inverses. Every source, Cartier, norm, character and
annular multiplier remains in these endpoint maps. The physical
return uses the inverse endpoint. This theorem assigns no new return
class. Its source domain is exactly \(\mathcal V\).
\end{theorem}
\begin{proof}[Supplied source argument; continuation review in progress]

% Expanded verified cp36 input: notes/part3_e6_intrinsic_base_proof.tex
\par\bigskip\noindent\textbf{Complete supporting proof 1: EB.}\par




\par

\medskip\noindent\textbf{The intrinsic elliptic quotient on the full noncollision E6 source locus}

\par\nobreak

For an arbitrary endpoint isomorphism we construct its induced
elliptic quotient before applying the Weierstrass invariant to fix the
source base. This treats the finite branch-permutation candidates left by
the auxiliary centroid polynomials in EF.

\par

\medskip\noindent\textbf{EB.1. The exact source and completed objects}

\par\nobreak

Retain
\[
 E_r:\quad y^2=x^3+rx^2+tx+t^2,\qquad G=(0,t),\qquad
 \mathcal V=\{r:r(4r-1)(3r-1)\ne0\}.
 \tag{EB.1}
\]
Choose either simple I1 as the selected one, with its original height
\(q=t-t_c\), complete branch and fan marking. Use either
\((a,b;m)=(6,1;3l+2)\), \(l\ge0\), or
\((a,b;m)=(2,3;3l)\), \(l\ge1\), either norm \(n=\pm1\),
and the prescribed finite character and full K96 reference.

The original source has precisely IV*, I2 and two simple I1 fibres
throughout EB.1, by the exact discriminant in EF.8. The root,
point-Cartier and annular maps EX.1-\/-19 and EF.28-\/-38 use
\(s=1/4-r\ne0\); the finite source charts and their full line
modules remain through every point of EB.1. A compact path in the
marked simple-node cover from the concrete EF.5 source to a given
object has a common smaller strict scalar contraction. All actual
source/root/line maps, the prescribed selected fan and the reference
continue on this path. The F46 and F35 integral source/flop topology
maps therefore give the same sphere, with \(H^2=H^4=0\).
The scalar can then be moved radially inside its contraction domain
to the specified value using the same complete cap square. This
gives the topology on EB.1 without any condition on \(P_r,Q_r\).

The horizontal-divisor and transverse \(t=u^e\) Laurent argument
EF.2 now applies literally. It proves
\(\mathbb C(X)=\mathbb C(t)\), with its specified pullback.
Every holomorphic isomorphism between two completed endpoints of the fixed
source thus descends to an actual Möbius base map \(h\).

\par

\medskip\noindent\textbf{EB.2. The full period module forces the original elliptic quotient}

\par\nobreak

On the regular base keep the full actual lattice of EF.11,
\[
 \gamma=(-bg,\ell),\qquad e_1=(1,0),\qquad
 e_2=(\tau,ag),\qquad\delta=(0,1),
 \tag{EB.2}
\]
and the complete monodromy EF.18-\/-19, including its central cochain.
The nonzero Abel-class proof EF.4 uses the nontrivial I1 monodromy
and the actual infinity logarithm of \(3G\), whose original
additive value is \(-s/2\). Both remain exactly as written on EB.1.
In particular the full invariant lattice line and invariant dual line
are intrinsically
\[
       \mathbb Z\delta,\qquad\mathbb Z\psi,
       \quad\psi(\gamma)=1,\quad\psi(e_1)=\psi(e_2)=\psi(\delta)=0.
 \tag{EB.3}
\]
The calculation establishing these lines does not assume that an
endpoint fixes the base: it is a computation of the invariant
subgroups of the local system itself. An isomorphism covering
\(h\) carries them to the corresponding subgroups of the pulled
back local system. Primitivity therefore gives
\(F_*\delta=\pm\delta\) and \(F^*\psi=\pm\psi\).

The subgroup \(\ker\psi\) gives the actual cyclic cover of a
regular torus fibre,
\[
 \widetilde T_t=
 \mathbb C^2/\langle e_1,e_2,\delta\rangle
           \simeq M_a^\times|_{E_{r,t}}.
 \tag{EB.4}
\]
The displayed identification is precisely the point-theta arrow
\(V_{\rm pt}=f_{P^a}(Z)v_M\), including its full constant factor
and the inverse Cartier-module maps EF.10-\/-17. A local lift of
\(F\) between these covers exists because it carries this exact
subgroup, and it is affine in the displayed universal-cover
coordinates: its periodic derivative is constant on the compact
torus fibre. Since the derivative takes the real vector
\(\delta=(0,1)\) to \(\pm(0,1)\) and is complex linear, it
takes its full complex span to itself. Its first coordinate is
therefore exactly
\[
                       z'=A(t)z+z_0(t),\qquad A(t)\ne0.
 \tag{EB.5}
\]
The two elliptic periods and the primitive \(\delta\) are carried
bijectively inside the actual kernel lattice. Quotienting that
lattice by \(\mathbb Z\delta\) gives
\[
          A(t)\langle1,\tau(t)\rangle
                =\langle1,\tau(h(t))\rangle.
 \tag{EB.6}
\]
Thus EB.5 descends to a degree-one affine holomorphic isomorphism
of the original elliptic source fibres. In multiplicative terms the
map first lifts to the genuine \(M_a^\times\), then preserves its
physical \(\mathbb C^*\) fibres with exponent \(\pm1\), and
then descends by the actual source projection. Its inverse follows
the same three arrows. The full line constants and the central
integer cochain remain in the lifted second coordinate; this
quotient assertion does not delete them from the completed map.

\par

\medskip\noindent\textbf{EB.3. The exact source invariant fixes three actual base points}

\par\nobreak

The Weierstrass coefficients give
\[
 c_4=16(r^2-3t),\quad
 \Delta=-16t^2[27t^2+(4-18r)t+r^2(4r-1)],
\]
\[
 j_r(t)=\frac{-256(r^2-3t)^3}
 {t^2[27t^2+(4-18r)t+r^2(4r-1)]}.
 \tag{EB.7}
\]
Because EB.5-\/-6 are actual elliptic isomorphisms, their source
invariant satisfies the literal equality
\(j_r(h(t))=j_r(t)\). Every relevant valuation is exact on EB.1.
The pole at zero has order two; the other two finite poles have
order one. The numerator has its triple zero at
\(t_z=r^2/3\), and at that point its remaining denominator factor is
\[
 27t_z^2+(4-18r)t_z+r^2(4r-1)
                  =\frac{r^2(3r-1)^2}{3}\ne0.
 \tag{EB.8}
\]
The unique other zero is infinity, of order one; in its exact
coordinate \(q_\infty=1/t\),
\(j_r(1/q_\infty)=256q_\infty+O(q_\infty^2)\).
Hence every base map preserving EB.7 fixes the unique double pole
zero, the unique simple zero infinity, and the unique triple zero
\(r^2/3\). These three points are distinct. The first two make the
Möbius map \(h(t)=ct\), and the third makes \(c=1\).

This proves pointwise preservation of the original base on all of
EB.1, including every candidate permutation of the two marked
simple-node branches. It is derived from the actual cyclic-cover
and elliptic-source maps, before any invariant is compared.

\par

\medskip\noindent\textbf{EB.4. Completion of the fixed-source group and endpoint scope}

\par\nobreak

With the base now fixed, the full centralizer and real-rank proof
EF.4 apply with all their cochains. The completed fixed-curve
obstruction EF.5 excludes its negative automorphisms. The entire
finite-source, selected-fan, integral source-generator, Miller and
root divisor proof EF.6 applies verbatim: its hypotheses are the
source fibre types of EB.1 and \(s\ne0\), including the separately
checked \(s=1\) point/O specialization. Both exact cap and annular
arrows EF.7 retain every norm sign and finite character.
Consequently
\[
              \operatorname{Aut}_{\rm hol}(X)
                    =\mathbb C^*_{\rm physical}
 \tag{EB.9}
\]
on EB.1, with the same actual injection and inverse-endpoint isotopy.

The scalar endpoint construction SEF.1-\/-4 uses those same source,
Cartier, root and reference arrows. Its integer winding restriction
uses precisely EB.3 and the full period intertwiner. The NE
selected vertex maps, genuine root cover, and retained full
annular integer \(n+m\nu\) are unchanged. Thus the complete
endpoint maps and winding groups SEF/NE.11 extend to EB.1 by these
literal identifications. The points \(r=0,1/4,1/3\) have different
actual source degenerations; no endpoint at those sources is
assigned by this note.

% End expanded verified cp36 input


\par\bigskip\noindent\textbf{Complete supporting proof 2: EBR.}\par


% Expanded verified cp36 input: notes/part3_e6_intrinsic_base_ebr_proof.tex
\par

\medskip\noindent\textbf{Root proof of the intrinsic base and the full noncollision endpoint locus}

\par\nobreak

We use the original source
\[
 E_r:\ y^2=x^3+rx^2+tx+t^2,\qquad G=(0,t),\qquad
 s=\tfrac14-r,\qquad
 \mathcal V=\mathbb C\setminus\{0,\tfrac14,\tfrac13\}.
 \tag{EBR.1}
\]
The completed threefold, its point line, its selected simple node,
its integral basis fan, and its entire annular reference are the
objects of EF.1, with either
\((a,b;m)=(2,3;3l)\), \(l\geq1\), or
\((a,b;m)=(6,1;3l+2)\), \(l\geq0\), and \(n=\pm1\).
Every endpoint in this proof is a holomorphic isomorphism between
the specified completed endpoints.

\par

\medskip\noindent\textbf{EBR.1. The marked parameter cover and the actual topology comparison}

\par\nobreak

Write
\[
 F_r(t)=27t^2+(4-18r)t+r^2(4r-1).
\]
The cover which chooses the selected simple node is exactly
\[
 w\longmapsto
 \left(r=\frac{1-w^2}{3},
 \quad t_c=\frac{(w-1)^2(2w+1)}{27}\right),
 \qquad
 w\in\mathbb C\setminus\{0,1,-1,\tfrac12,-\tfrac12\}.
 \tag{EBR.2}
\]
Its inverse on \(F_r(t_c)=0\), \(r\in\mathcal V\), is
\[
 w=\frac{54t_c-18r+4}{4(1-3r)}.
 \tag{EBR.3}
\]
Substitution proves both composites are the identity. In particular
the other selected-node choice is \(w\mapsto-w\), with the same
source \(r\). The displayed punctured plane is path connected:
a polygonal path between two points can be perturbed in finitely
many disjoint small discs to avoid its five excluded points.
This proves the precise connectedness assertion used in EB.1.

Along such a compact path the two simple nodes remain distinct
and nonzero; their exact discriminant is \(16(1-3r)^3\).
At either node the parameter-dependent analytic Morse coordinates
give the same original height \(q=t-t_c\). The I2 coordinate
charts retain \(s\ne0\). The infinity root has equation
\[
 Y^2=X^3+r z^2X^2+zX+1,
 \quad q_\infty=z^3,\quad x=z^{-2}X,\quad y=z^{-3}Y,
 \tag{EBR.4}
\]
and its central curve and deck weights are unchanged. The higher
good and reduced roots, their point/O graph maps, and their complete
chord lifts are precisely EF.32-\/-38, with the same varying \(r\).
All finite characters and the chosen reference cochain are continued
through those actual charts.

The strict contraction inequalities are open on each compact norm
shell. A finite cover of the compact path gives one smaller nonzero
scalar for every chart in that path. The source, root, selected-fan
and annular squares then glue a proper smooth family of the completed
spaces over the path. An explicit smooth comparison is obtained by
lifting the path tangent field in these charts and patching the lifts
with a partition of unity on the common compact shell. On each glued
boundary the lifts are transported by its actual gluing differential;
thus the patched field descends to the completed family. Its flow
exists throughout the compact parameter interval, with inverse given
by the reverse field, and identifies its endpoint fibres. The
integral F35/F46 maps give the sphere at the chosen EF.5 starting
source and every prescribed integral fan. Transport the prescribed
fan marking backwards along the path to that starting source; the
same flow then applies to that marking. Radial transport within the
strict contraction domain reaches the required scalar. Consequently
the actual completed objects throughout EBR.1 retain
\(H^2(X;\mathbb Z)=H^4(X;\mathbb Z)=0\).

The horizontal-divisor argument EF.7 uses exactly this vanishing
and a flat Kähler form on a regular two-torus fibre. Its equality
between the divisor class pairing and positive divisor area excludes
horizontal polar divisors. The transverse map at every special
component is \(u\mapsto u^e=t-t_0\); invariance under its \(e\)
branches makes the descended Laurent exponents integral. Thus the
specified pullback identifies \(\mathbb C(X)\) with
\(\mathbb C(t)\). An endpoint isomorphism \(F\) therefore satisfies
\(\pi_{\rm target}F=h\pi_{\rm source}\) for an actual Möbius map
\(h\), including at the special fibres by holomorphic continuation.
The derivative of this equation and its inverse show that \(h\)
permutes the critical values and carries the regular locus onto
the regular locus.

\par

\medskip\noindent\textbf{EBR.2. The complete lattice gives the original elliptic morphism}

\par\nobreak

Retain the entire period module, including the source-dependent
suspension logarithm and the integer central cochain:
\[
 \gamma=(-bg,\ell),\quad e_1=(1,0),\quad
 e_2=(\tau,ag),\quad \delta=(0,1),\qquad
 M_\alpha=
 \begin{pmatrix}
 1&0&0\\ bv_\alpha&A_\alpha&0\\
 c_\alpha&a v_\alpha^tJA_\alpha&1
 \end{pmatrix}.
 \tag{EBR.5}
\]
The periods are mapped to the point-Cartier line by the full
\(V_{\rm pt}=f_{P^a}(Z)v_M\) arrow EF.10, with the constant
in EF.12 and the reciprocal module maps at both divisors.
No coefficient is replaced by a central value.

The nonzero rational Abel-class proof EF.4 applies at every
parameter in EBR.1. At the finite multiplicative fibres an invariant
period is the extending annular period. At infinity the actual
logarithm of \(3G'\) is \(-sz^2/2+O(z^5)\); division by \(z^2\)
gives the original extending Lie frame. A rational coboundary would
therefore give a global section of the Lie line \(\mathcal O(-1)\).
The explicit coefficient transition
\(u_0(q)=q^{-1}u_t(q^{-1})\) makes this section zero. The original
central additive coordinate of \(3G\) is \(-s/2\ne0\), so its
nonzero integer multiples cannot be the zero section. This proves
the required nonzero Abel class on the whole stated domain.
The infinity elliptic matrix has no invariant vector. In EBR.5 an
invariant vector with \(\gamma\)-coefficient \(z\) would satisfy
\((A_\alpha-I)u+bzv_\alpha=0\); hence \(z=0,u=0\).
The dual calculation, using the retained
\(a v_\alpha^tJA_\alpha\), gives the dual invariant line.
Primitivity of the displayed lattice generators yields exactly
\[
 \Lambda^{\pi_1}=\mathbb Z\delta,\qquad
 (\Lambda^*)^{\pi_1}=\mathbb Z\psi,\qquad
 \psi(\gamma)=1,\quad\psi(e_1)=\psi(e_2)=\psi(\delta)=0.
 \tag{EBR.6}
\]

The endpoint intertwines this local system with its pullback by
\(h\). The invariant subgroups in EBR.6 are preserved under this
intertwining, without assuming \(h=1\). Thus
\(F_*\delta=\epsilon\delta\) and \(F^*\psi=\epsilon'\psi\),
with \(\epsilon,\epsilon'\in\{1,-1\}\). In particular it carries
\[
 H_t=\ker\psi_t=\langle e_1,e_2,\delta\rangle_{\mathbb Z}
 \quad\hbox{bijectively to}\quad H_{h(t)}.
 \tag{EBR.7}
\]
The actual corresponding cyclic cover is
\(\mathbb C^2/H_t=M_a^\times|_{E_{r,t}}\), by the point-theta
arrow just specified.

Lift the endpoint on a regular fibre to its universal cover.
The derivative of that lift is periodic for the full rank-four
lattice; each matrix entry therefore descends to a holomorphic
function on a compact complex torus and is constant. The lift is
affine. Complex linearity and \(F_*\delta=\epsilon\delta\)
give its exact form
\[
 (z,w)\longmapsto
 (A z+z_0,\ C z+\epsilon w+w_0),\qquad A\ne0.
 \tag{EBR.8}
\]
Its inverse is
\[
 (z',w')\longmapsto
 \left(\frac{z'-z_0}{A},
 \epsilon\left(w'-w_0-C\frac{z'-z_0}{A}\right)\right).
 \tag{EBR.9}
\]
The full second-coordinate constants remain in these arrows.
On the physical nonzero-line fibre the multiplicative map is
\(V'\!=e^{2\pi i(Cz+w_0)}V^\epsilon\), interpreted in its
actual source and target point-theta frames.

The lattice map on EBR.7 has the exact integral form
\[
 \begin{pmatrix}p&q&0\\ r_0&s_0&0\\ u&v&\epsilon\end{pmatrix},
 \qquad
 A=p+r_0\tau(h(t)),\quad
 A\tau(t)=q+s_0\tau(h(t)).
 \tag{EBR.10}
\]
It is unimodular. The first-coordinate complex map preserves real
orientation, so \(ps_0-qr_0=1\). These equalities prove
\[
 A\langle1,\tau(t)\rangle_{\mathbb Z}
 =\langle1,\tau(h(t))\rangle_{\mathbb Z}.
 \tag{EBR.11}
\]
Consequently \([z]\mapsto[Az+z_0]\) and its displayed inverse
are degree-one holomorphic maps between the original elliptic
source fibres. They are induced by EBR.8-\/-9 through the literal
cyclic cover and the physical \(\mathbb C^*\) projection.
The central integers \(u,v,c_\alpha\) remain in the lifted line
map. This supplies both directions of the geometric comparison
needed to use the elliptic invariant.

\par

\medskip\noindent\textbf{EBR.3. The exact invariant fixes the entire original base}

\par\nobreak

Translate the cubic variable by \(x=X-r/3\). Its exact short
Weierstrass coefficients are
\[
 p=t-r^2/3,\qquad q=t^2-rt/3+2r^3/27.
 \tag{EBR.12}
\]
The formula \(1728\,4p^3/(4p^3+27q^2)\) therefore gives
\[
 j_r(t)=
 -\frac{256(r^2-3t)^3}{t^2F_r(t)}.
 \tag{EBR.13}
\]
The translated-cubic computation is evaluated symbolically in \path{code/current_audits/audit_corrections/code/independent_IV_endpoint_checks.py}. EBR.8-\/-11 give the actual
elliptic isomorphisms, hence \(j_r\circ h=j_r\).

At zero the Laurent coefficient is
\[
 \lim_{t\to0}t^2j_r(t)=-\frac{256r^4}{4r-1}\ne0.
 \tag{EBR.14}
\]
The other poles, the two distinct roots of \(F_r\), are simple.
Indeed the only finite numerator zero is \(t_z=r^2/3\), and
\[
 F_r(t_z)=r^2(3r-1)^2/3,\qquad
 t_z^2F_r(t_z)=r^6(3r-1)^2/27,
 \tag{EBR.15}
\]
both nonzero on EBR.1. Thus \(t_z\) is the unique triple zero.
At infinity the original coordinate gives
\[
 j_r(1/q_\infty)=256q_\infty+O(q_\infty^2),
 \tag{EBR.16}
\]
so infinity is the unique simple zero. Pullback by a Möbius map
preserves each local order, by substitution of its invertible
local coordinate. Therefore \(h\) fixes zero, infinity and
\(r^2/3\). The first two equalities give \(h(t)=ct\); the
third, with \(r\ne0\), gives \(c=1\).

This argument includes the actual auxiliary-polynomial zeros
excluded in EF.3. As a concrete branch-permutation candidate,
\(r=2/9\) has \(F_r(t)=27t^2-4/729\); the map \(t\mapsto-t\)
permutes its two simple poles but moves the unique triple zero
\(4/243\). Substitution into EBR.13 also gives an explicit
nonzero rational difference. The geometric quotient above,
followed by these exact orders, excludes that candidate on the
completed objects.

\par

\medskip\noindent\textbf{EBR.4. Every remaining endpoint map retains its full domain}

\par\nobreak

After \(h=1\), the entire EF.18 centralizer computation applies
to EBR.5, including \(c_\alpha\). It gives
\(eI+c\delta\psi\); complex linearity removes its real rank-one
term and gives the two signs. For an automorphism of one completed
object, the negative sign has sixteen transverse fixed points in
every regular fibre. Their proper analytic curve closure has
positive degree over the original base, whereas pairing with
\(\pi^*[\omega_{\mathbb P^1}]=0\) gives zero. The full completed
fixed-set argument EF.5 thus excludes this sign throughout EBR.1.

For the positive sign, the same actual cyclic-cover lifts and
selected-vertex unit arrows EF.24-\/-25 give a global original
source section. Its component labels remain those carried by
the finite I2 and infinity root charts. The full ten-class Gram
matrix EF.26a has determinant \(-1\); the two disjoint section
charts and the E6/A1 correction columns give height \(1/6\).
Thus the source section is \(jG\), \(j\in\mathbb Z\), on all
of EBR.1. Its actual line lift is \(H(t)U_G^j\).

No discarded auxiliary polynomial enters its three divisor
orders. At the I2 the M2 full coefficient is
\(t(\beta-1-\alpha/2)/(-s+t\alpha)\). The M6 full coefficient,
after the exact original \(t^3\) factor, is
\(8((1-s)w-s)/s^3\). Its separately computed \(s=1\)
point/O specialization is \(-8\), with the original two
Cartier equations retained. The root maps EF.32-\/-35 have
generic inverse orders \(2m/b\) and all their point/O correction
monomials. At the selected node the actual \(B^k\) reindex
gives its integer order \(-k\).
The full orders of \(H\) are therefore \(3j/b,-2j/b,-k\)
at zero, infinity and the selected node. Its divisor degree
gives \(j=bk\) and
\[
 H=c(t^3/q)^k,\qquad H U_G^{bk}=cB_{\rm sel}^{-k}.
 \tag{EBR.17}
\]
On the quotient this is precisely multiplication by
\(c\lambda^k\). It commutes with every actual cap and annular
arrow EF.34-\/-38. The additive value \(-s/2\) also proves the
physical scalar injection has trivial kernel. Thus
\(\operatorname{Aut}_{\rm hol}(X)=\mathbb C^*_{\rm physical}\)
on the entire stated locus.

For scalar endpoints, the full ordinary intertwiner and its
two complex-independent vectors \(e_1,\delta\) give SEF.21:
the two original scalars are equal and the path winding \(J\)
is integral. The root-side comparison is the full
\(z^{nJ}\) in SEF.5. The third divisor equation becomes
\[
 m\,\operatorname{ord}_\infty H+2mj/b+nJ=0.
 \tag{EBR.18}
\]
Together with the other two actual orders it gives
\(j=bk+bnJ/m\), hence \(m\mid bJ\).
For M2 the full primitive forward maps are
\[
 (R^{-l}tU_G)^n\quad\hbox{on the ordinary piece},\qquad
 (R^{-l}U'_G)^n\quad\hbox{on the root piece},
 \quad R^{3l}=\lambda .
 \tag{EBR.19}
\]
The complete SEF.14 Cartier comparison contains \(z^{2l}\);
its product with \(z^{nl}\) and \(t^n=z^{-3ln}\) is one.
This is the actual annular endpoint square, for either signed
inverse and every finite character. Its three chronological
composites are \(q^n\) and \(h^n\), with
\(h=1-t_cz^m\), as in SEF.20. For M6 the primitive maps
are these \(q^n/h^n\) arrows themselves.

The negative scalar-endpoint obstruction uses the actual
ordinary and root covers of NE.2-\/-3. The complete reference
monodromy on one root circle is \(B_\lambda^{n+m\nu}\),
with its retained integer \(\nu\). The two single-valued
negative lifts would give
\(B_\lambda^{2(n+m\nu)}=1\); strict contraction and
\(n=\pm1,m\ge2\) exclude this equality. The proof applies
to the quadratic M6 cap as well. All root, selected-vertex
and Cartier charts used there have just been retained over
EBR.1. Hence the full scalar winding groups are
\(l\mathbb Z\) for M2 and \(m\mathbb Z\) for M6, with the
explicit primitive arrows above and physical scalar ambiguity.
No scalar return class is computed or divided by this endpoint
classification. The three removed source parameters are outside
the scope of every assertion in this proof.

% End expanded verified cp36 input


\par\bigskip\noindent\textbf{Complete supporting proof 3: EBP.}\par


% Expanded verified cp36 input: notes/part3_e6_intrinsic_base_ebp_proof.tex
\par

\medskip\noindent\textbf{Intrinsic elliptic quotient and endpoint domain}

\par\nobreak

This peer proof concerns the complete EB and EBR constructions and their
use of EF, SEF and NE. The source is
\[
 y^2=x^3+rx^2+tx+t^2,\qquad G=(0,t),\qquad
 s=\tfrac14-r,\qquad r(4r-1)(3r-1)\ne0.
 \tag{EBP.1}
\]
The point line, selected simple node, ordered fan, scalar, norm sign,
finite character and complete annular reference are the original ones.
The cap data are \((a,b;m)=(6,1;3l+2)\), \(l\ge0\), and
\((2,3;3l)\), \(l\ge1\), with \(n=\pm1\). Every endpoint in
this argument is a holomorphic isomorphism of the specified completed
endpoints. Its inverse is retained in each quotient construction below.

\par

\medskip\noindent\textbf{EBP.1. The actual marked cover and topology transport}

\par\nobreak

Put \(F_r(t)=27t^2+(4-18r)t+r^2(4r-1)\). The selected-node cover is
\[
 r=(1-w^2)/3,\qquad t_c=(w-1)^2(2w+1)/27,
 \qquad
 w=\frac{54t_c-18r+4}{4(1-3r)}.
 \tag{EBP.2}
\]
Direct substitution gives both inverse composites, using \(F_r(t_c)=0\)
for the second. The original excluded product pulls back exactly to
\[
 -\frac{w^2(w^2-1)(4w^2-1)}9.
 \tag{EBP.3}
\]
Thus the parameter cover is the plane with precisely
\(0,\pm1,\pm\tfrac12\) removed. Its polygonal paths, detoured in
disjoint discs around these five points, join the concrete EF.5 source
to either marked branch of every object in EBP.1. The involution
\(w\mapsto-w\) gives the other root of the same quadratic; their sum
and product are the exact coefficients of \(F_r\).

Along a compact such path, the two selected candidates are distinct and
nonzero, since the discriminant is \(16(1-3r)^3\) and the constant
coefficient is \(r^2(4r-1)\). Analytic Morse charts keep the actual
height \(q=t-t_c\). The I2 charts use the unit \(s\). At infinity the
good cubic source root is literally
\[
 q_\infty=z^3,\qquad x=z^{-2}X,
 \quad y=z^{-3}Y,\qquad
 Y^2=X^3+rz^2X^2+zX+1.
 \tag{EBP.4}
\]
The higher roots are the full EF.32-\/-38 maps reviewed in EBP.6 below.
Their characters and reference integers are continued as part of those
maps. Transport the prescribed terminal fan backwards to the starting
object and then use these actual cone and period identifications.

A common smaller scalar gives strict contraction on a finite collection
of compact norm shells covering the path. In each source/root/fan chart,
lift the path vector field; on common boundary charts use the displayed
gluing derivative to identify the lifts. A partition of unity on the
compact glued family preserves the projection of the field to the path
tangent. Its flow and reverse flow identify the completed endpoint
fibres. This is the precise smooth comparison used with the original
F35/F46 integral topology maps. Radial scalar transport uses the same
full cap and reference square. Hence the resulting completed spaces
retain the sphere cohomology, including \(H^2=H^4=0\), throughout
EBP.1. The argument carries each prescribed discrete marking from the
corresponding initial object.

The horizontal-divisor pairing EF.7 now excludes every horizontal polar
divisor on a regular compact two-torus fibre. Removal across its remaining
codimension-two indeterminacy makes a meromorphic function constant on
that fibre. At each special component the actual transverse map
\(u\mapsto u^e=t-t_0\), after its extending unit root, forces the
Laurent exponents to be multiples of \(e\). The descended function is
therefore meromorphic at the special base value. This proves the precise
pullback equality \(\mathbb C(X)=\mathbb C(t)\). An endpoint and its
inverse induce an automorphism of this rational function field and give
\(\pi_{\rm target}F=h\pi_{\rm source}\) for a Möbius map. Taking
derivatives in that equality and its inverse carries the regular locus
onto the regular locus. No assertion that \(h\) fixes the base has yet
been used.

\par

\medskip\noindent\textbf{EBP.2. The two primitive invariants before base rigidity}

\par\nobreak

In the actual point-theta coordinates, retain
\[
 \gamma=(-bg,\ell),\quad e_1=(1,0),\quad e_2=(\tau,ag),
 \quad\delta=(0,1),\qquad V_{\rm pt}=f_{P^a}(Z)v_M.
 \tag{EBP.5}
\]
The full constant relating this coefficient to the tensor coefficient,
the two Cartier generators and their inverse maps are EF.10-\/-17.
In particular \(\ell\) is the logarithm of the entire suspension
coefficient. It retains the scalar, original selected height, Miller
function and point/O factors.

The monodromy, with \(J=\left(\begin{smallmatrix}0&-1\\1&0\end{smallmatrix}\right)\), is
\[
 M_\alpha=\begin{pmatrix}
 1&0&0\\ bv_\alpha&A_\alpha&0\\
 c_\alpha&a v_\alpha^tJA_\alpha&1
 \end{pmatrix}.
 \tag{EBP.6}
\]
Multiplication gives the full central term
\(c_{\alpha\beta}=c_\alpha+c_\beta+
abv_\alpha^tJA_\alpha v_\beta\), as well as
\(v_{\alpha\beta}=v_\alpha+A_\alpha v_\beta\). The peer
calculator verifies this identity with symbolic central coefficients.

Here is the domain of the nonzero Abel-class argument. A rational
coboundary, after clearing denominators and multiplying by six, gives
a single-valued Lie logarithm of an integer multiple of \(G\) lying in
every identity component. At the finite multiplicative fibres the local
Tate coordinate is an extending nonzero unit. The difference between
its logarithm and the global logarithm is an invariant annular period,
which extends in the same Lie frame. At infinity the actual chords give
\[
 x(3G)=4t^2/s^2+2t/s,\qquad
 y(3G)=8t^3/s^3+6t^2/s^2+t/s-t.
 \tag{EBP.7}
\]
On EBP.4 the formal parameter has leading term \(-sz^2/2\).
The original minimal differential is \(z^{-2}\omega_{\rm good}\),
whereas the affine differential is \(z\omega_{\rm good}\).
Dividing that root logarithm by \(z^2\) gives the extending original
Lie coefficient. The infinity elliptic monodromy has polynomial
\(x^2+x+1\), so it supplies no invariant-period ambiguity there.

The full Lie coefficient transition is
\(u_0(q_\infty)=q_\infty^{-1}u_t(q_\infty^{-1})\). Its global
holomorphic section is zero, because the entire function \(u_t\)
decays at infinity. But the actual original central point has
\((x_0,y_0)=(4/s^2,8/s^3)\), so its additive coordinate
\(-x_0/y_0=-s/2\) and all its nonzero integer multiples are nonzero.
The differential identity on \((x_0,y_0)=(u^2,u^3)\) is
\(dx_0/(2y_0)=d(-1/u)\). This proves the nonzero rational Abel class
for every parameter in EBP.1, using only \(s\ne0\) and the retained
fibre types.

An invariant vector with first coefficient \(z\) satisfies
\((A_\alpha-I)u=-bzv_\alpha\). The nonzero Abel class gives
\(z=0\); the infinity matrix gives \(u=0\). Its last coefficient
is free. For the dual equation write the middle row as \(w^tJ\).
The exact identity
\[
 w^tJ(A-I)=(w-Aw)^tJA
 \tag{EBP.8}
\]
turns the dual invariant equation into the same Abel coboundary
equation. Its final coefficient is zero and its middle row is zero.
Consequently the full integral invariant groups are precisely
\[
 \Lambda^{\pi_1}=\mathbb Z\delta,
 \qquad (\Lambda^*)^{\pi_1}=\mathbb Z\psi,
 \qquad\psi(\gamma)=1,
 \quad\psi(e_1)=\psi(e_2)=\psi(\delta)=0.
 \tag{EBP.9}
\]
The displayed basis proves primitivity. The equations include the entire
central cochain; it has not been assigned a different value.
An endpoint covering \(h\) intertwines this local system with its
pullback. It carries these intrinsic invariant groups bijectively,
so it sends \(\delta\) and \(\psi\) to their signed primitive
generators. This statement precedes the determination of \(h\).

\par

\medskip\noindent\textbf{EBP.3. The actual cyclic cover, affine arrow and elliptic inverse}

\par\nobreak

The subgroup
\(H=\ker\psi=\langle e_1,e_2,\delta\rangle\) defines the
literal cyclic cover \(\mathbb C^2/H=M_a^\times|_{E_t}\), through
the full point-theta and inverse Cartier-module arrows in EBP.5.
The endpoint carries \(H\) bijectively to the corresponding target
subgroup. Lift it to the universal covers. Its derivative is periodic
for the full rank-four lattice; every entry descends to a holomorphic
function on the compact torus and is constant. The lift is affine.
Complex linearity and the primitive physical period give its exact form
\[
 (z,w)\longmapsto(Az+z_0,Cz+\epsilon w+w_0),\qquad
 \epsilon\in\{1,-1\},\quad A\ne0,
 \tag{EBP.10}
\]
with inverse
\[
 (z',w')\longmapsto
 \left((z'-z_0)/A,
 \epsilon\{w'-w_0-C(z'-z_0)/A\}\right).
 \tag{EBP.11}
\]
In multiplicative point coefficients the full line arrow is
\(V'=\exp(2\pi i(Cz+w_0))V^\epsilon\). Its target and source
Cartier frames are the ones already specified; their constant factors
and integral transitions remain in \(C,w_0\) and the period rows.

The map on the kernel lattice is an integral matrix
\[
 \begin{pmatrix}p&q&0\\r_0&s_0&0\\u&v&\epsilon\end{pmatrix},
 \quad A=p+r_0\tau',\quad A\tau=q+s_0\tau'.
 \tag{EBP.12}
\]
It is unimodular. The area identity
\[
 \operatorname{Im}\bigl(\overline{p+r_0\tau'}(q+s_0\tau')\bigr)
 =(ps_0-qr_0)\operatorname{Im}\tau'
 =|A|^2\operatorname{Im}\tau
 \tag{EBP.13}
\]
gives \(ps_0-qr_0=1\). Thus the quotient period lattices agree
exactly. The maps \([z]\mapsto[Az+z_0]\) and EBP.11\textquotesingle s first
coordinate are inverse holomorphic isomorphisms of the original
elliptic source fibres. They are obtained through the specified
cyclic-cover and physical-fibre projections. The integers \(u,v\)
and the full central cochain remain in the line map, where they belong.

\par

\medskip\noindent\textbf{EBP.4. Exact base rigidity on the whole stated locus}

\par\nobreak

Translating \(x=X-r/3\) gives the short cubic coefficients
\(p=t-r^2/3\), \(q=t^2-rt/3+2r^3/27\). Direct computation gives
\[
 c_4=16(r^2-3t),\qquad\Delta=-16t^2F_r(t),\qquad
 j_r(t)=-\frac{256(r^2-3t)^3}{t^2F_r(t)}.
 \tag{EBP.14}
\]
The actual source isomorphisms EBP.10-\/-13 prove \(j_r\circ h=j_r\).
Its double-pole coefficient at zero is \(-256r^4/(4r-1)\).
The other two poles are simple. The unique finite zero is the triple
zero \(t_z=r^2/3\), since
\[
 F_r(t_z)=r^2(3r-1)^2/3,\qquad
 t_z^2F_r(t_z)=r^6(3r-1)^2/27.
 \tag{EBP.15}
\]
Infinity is the unique simple zero, with coefficient
\(j_r(1/q_\infty)=256q_\infty+O(q_\infty^2)\).
All these nonzero coefficients are valid on EBP.1. An invertible local
coordinate preserves the order of a zero or pole, so \(h\) fixes
zero, infinity and \(r^2/3\). The first two give \(h(t)=ct\), and
the third gives \(c=1\).

This includes the auxiliary centroid-polynomial zeros from EF.3.
For example, at \(r=2/9\), \(F_r(t)=27t^2-4/729\); its pole
permutation \(t\mapsto-t\) moves \(t_z=4/243\). The exact
rational functions \(j_r(-t)\) and \(j_r(t)\) differ. The source
isomorphism was established before making this invariant comparison.

\par

\medskip\noindent\textbf{EBP.5. The remaining automorphism hypotheses and source generator}

\par\nobreak

After the base is fixed, the full centralizer calculation is
\(eI+c\delta\psi\). Its commutator with EBP.6 vanishes for every
retained central coefficient. The real rank-one remainder is eliminated
by complex linearity, and unimodularity gives \(e=\pm1\). For an
automorphism of a completed object, the negative sign gives sixteen
transverse fixed points per regular two-torus. Their closed analytic
equalizer in the compact completed space has a horizontal compact curve
component. Its positive base degree contradicts the already transported
\(H^2=H^4=0\). This is the full completed fixed-curve argument EF.5.

For the positive sign, the ordinary cyclic-cover extension, genuine
root pullback, and selected vertex unit maps EF.24-\/-25 give the actual
global source section. All vanishing cycles lie in \(\ker\psi\);
the finite cap uses \(A^m=B^n\) and strict contraction to remove
the commutator integer. The lift torsor over the one-puncture orbifold
base has no homomorphism from its finite cyclic fundamental group to
the suspension deck group. These are the explicit extension steps,
with both signs and all reference integers retained.

The rational source chart is
\(U=y-t-x/2\), \(V=y+t+x/2\), \(UV=x^2(x-s)\), with rational
inverse \(t=(V-U-x)/2\), \(y=(V+U)/2\). The I2 tangent cone
\(UV=-sx^2\) and the two distinct infinity sections retain the
component labels used in EF.26a. The peer calculator reconstructs the
full ten-class Gram matrix from those labels and obtains determinant
\(-1\). The E6 correction is \(4/3\), the A1 correction is
\(1/2\), and the height of \(G\) is \(1/6\). Hence the actual
global section is an integral multiple \(jG\), with the same
primitive generator and no omitted section torsion.

The full endpoint lift is \(H(t)U_G^j\). For M2 the exact finite chart
gives \(B_{2,1}/t=2v/(s-vh)\). For M6 its coefficient after the
actual \(t^3\) factor is \(8((1-s)w-s)/s^3\), which is not
identically zero for any \(s\ne0\). The special value \(s=1\)
is checked from the full chord before taking a limit. The independent
calculation retains the actual tangent at \(3G\), the doubled point,
and both translated quotients, and obtains
\[
 \frac{f_3(p-G)}{h f_3(p)}\longrightarrow
 \frac{w^3}{(w-1)(w^2-w+1)},\qquad
 \frac{f_{2,3G}(p-G)}{h f_{2,3G}(p)}\longrightarrow
 -\frac{(w-1)^2}{w^2-w+1}.
 \tag{EBP.16}
\]
Their complete product gives \(B_{6,1}/t^3\to-8\), and
\(\tau(6G)/t\to-6\). The two actual point/O Cartier equations
are retained at that coincidence. The Hom divisor and its inverse
therefore extend through the resolved chart by the same normality
argument as EF.30g. This explicitly covers the only extra specialization
in that finite-source factor calculation.

The three original divisor orders of \(H\) are
\(3j/b,-2j/b,-k\) at zero, infinity and the selected node.
Their degree gives \(j=bk\) and the literal equality
\(H U_G^{bk}=cB_{\rm sel}^{-k}\). On the actual quotient this
is multiplication by \(c\lambda^k\). The central additive value
from EBP.7 makes the physical scalar injection faithful. Its full
source/root/fan extension and inverse are the EF.34-\/-38 maps.

\par

\medskip\noindent\textbf{EBP.6. Both root families and every downstream endpoint square}

\par\nobreak

The higher-root substitutions retain the following complete exponents:
\[
\begin{array}{c|c|c}
(a,b;m)&(q_\infty,x_0,y_0)&\text{root equation}\\ \hline
(6,1;3l+2)&(z^{3l+2},z^{4l+2}X,z^{6l+3}Y)
 &Y^2=X^3+rz^{2l+2}X^2+z^{l+2}X+z^2\\
(2,3;3l)&(z^{3l},z^{4l}X,z^{6l}Y)
 &Y^2=X^3+rz^{2l}X^2+z^lX+1.
\end{array}
 \tag{EBP.17}
\]
The reduced source division and all graph rays are EF.32; the good
division and both point/O correction sums are EF.33. The parameter
\(r\) is an ordinary holomorphic coefficient in these equations.
The formal sections and normalized chord lifts use \(s\ne0\) and
the same original central points. Their full norm unit has central
value four. The prescribed root, character \(\zeta_m^d\), signed
norm and entire reference give the same \(A_\lambda\), \(B_\lambda\)
and sector maps EF.35-\/-38. The nonzero selected factor is exactly
\(h=1-t_cz^m=q_\infty q\), with value one at the root centre.
No centroid polynomial occurs in any of these maps.

For SEF the actual finite map is \(S=tU_G\). On the good root,
\(f_{2,\rm old}=z^{-l}\phi\),
\(U_{G,\rm old}=z^{2l}U'_G\), with
\(\phi=(Y-1-z^lX/2)/(X-sz^{2l})\).
Its central product is \(-X\) in the retained meromorphic point
frames, with exactly their Cartier divisor. The actual graph correction
\(E_P-E_O\) remains in both directions of that equality.
The full deck calculation gives
\(U'_G D=\zeta_m^{-2l}D U'_G\). Consequently the primitive
forward maps are
\[
 (R^{-l}tU_G)^n\quad\hbox{and}\quad(R^{-l}U'_G)^n,
 \tag{EBP.18}
\]
with the directed finite square SEF.18. For \(n=-1\) the
coefficients are evaluated at the actual inverse source point.

The full reference restoration SEF.3 retains every sector integer.
At integral winding \(J\) its endpoint ratio is exactly
\(z^{nJ}\), since its additional scalar factor on an integral
suspension label is one. For the primitive good endpoint the three
annular powers are \(nl,-3nl,2nl\); their sum is zero. The three
chronological endpoint maps retain the successive root scalars and
compose to \(q^nB_\lambda^{-n}\) and \(h^nB_\lambda^{-n}\).
All finite characters commute through these squares and retain their
specified norm. This proves the primitive and iterated maps on the
entire domain EBP.1.

For an arbitrary endpoint the complete ordinary periods first force
the two original scalars to agree and their path winding to be an
integer \(J\). The three full divisor equations then give
\(j=bk+bnJ/m\), hence \(m\mid bJ\). With the explicitly
constructed primitive maps this yields \(l\mathbb Z\) for M2
and \(m\mathbb Z\) for M6.

The negative-sign exclusion uses the full NE construction. Its selected
inversion carries the two actual vertex charts by NE.4, preserving
their product \(q\), their full unit coefficients and their inverse.
Composition with that inversion supplies the single-valued ordinary
cyclic-cover lift across the selected puncture. The original source/O
section and point/O module frames extend the other ordinary lifts.
The genuine root fibre product similarly gives a single-valued root
lift fixing \(z\). Its full annular deck is
\(A_\lambda B_\lambda^\nu\), with the actual carried integer
\(\nu\). Its chronological \(m\)-fold product is
\(B_\lambda^{n+m\nu}\); the same integer occurs at the two
scalar endpoints by the explicit SEF restoration.
Transporting the negative lift around the root circle gives the
conjugate factor \(B_\lambda^{2(n+m\nu)}\). Single-valuedness
and strict contraction imply \(n+m\nu=0\), whereas this integer
is congruent to \(\pm1\) modulo \(m\ge2\). Thus the full NE
obstruction, including the quadratic reduced cap, retains its proof
on EBP.1. The actual winding maps and physical scalar ambiguity are
therefore exactly the stated ones.

\par

\medskip\noindent\textbf{EBP.7. Independent verification and scope}

\par\nobreak

The invariant calculation uses the original cubic and the displayed
marked-cover composites. The elliptic quotient and its inverse, the
primitive kernel and dual equations, the nonzero Lie class, and the
compact flow retain the original source, point line and annular reference.
The exact source and matrix identities are evaluated in \path{code/current_audits/audit_corrections/code/independent_IV_endpoint_checks.py}.

The completed independent run passes 2,396 exact checks, including 264
complete integral kernel matrices with both inverse products. The
symbolic identities cover the arbitrary central cochain and arbitrary
signed root exponent; the finite executions separately retain all
displayed coefficients, source inverses and endpoint words.

At the base-rigidity step we use EBP.2--4's actual elliptic quotient
and invariant. EBP.1 and the maps above retain the original source,
Cartier module, roots, periods, references and endpoint-isomorphism
hypotheses throughout the stated noncollision locus.

% End expanded verified cp36 input



\par\medskip\noindent\textbf{Source and endpoint inputs.}\par
The source/root extension and physical group are
\cref{night:P3K.109}; the scalar endpoints, full references and negative
lifts are \cref{night:P3K.110}. EBR.4 and EBP.5--6 give the source,
Cartier, root, finite-character, period and endpoint maps on the
noncollision locus stated here.

The program \path{code/current_audits/audit_corrections/code/independent_IV_endpoint_checks.py} derives the original cubic invariant,
the marked-node parameter and its inverse, the exact zero and pole
coefficients, the full symbolic central-cochain product, the dual-invariant
equations and the ten-class determinant. It retains the literal specialized
$M_6$ Miller ratios and whole-unit affine inverse. The output
\path{code/current_audits/audit_corrections/independent_IV_endpoint_checks.json} gives these identities and the finite signed
endpoint checks. These are the algebraic inputs to the elliptic quotient,
Lie extension, proper closure and compact comparison proved above.


% End expanded verified cp36 input

\end{proof}


\subsection{Full holomorphic endpoints on the original type-III E6 source}
